\section{Algebra of limits and indeterminate forms}

Limit computations often use algebra. But algebra should be used with understanding. Some transformations preserve the behavior near a point; others may accidentally change the problem. The guiding question should always be: what happens to the expression as the input approaches the limiting value?

\subsection*{Basic algebra of limits}

If two limits exist,
\[
\lim_{x\to a}f(x)=L,
\qquad
\lim_{x\to a}g(x)=M,
\]
then we can combine them in predictable ways:
\[
\lim_{x\to a}(f(x)+g(x))=L+M,
\]
\[
\lim_{x\to a}(f(x)g(x))=LM,
\]
and, if \(M\ne0\),
\[
\lim_{x\to a}\frac{f(x)}{g(x)}=\frac{L}{M}.
\]

These rules are useful when the individual limits are known and the resulting expression is meaningful.

\begin{examplebox}
Compute
\[
\lim_{x\to2}(x^2+3x-1).
\]
Using direct substitution and algebra of limits,
\[
\lim_{x\to2}(x^2+3x-1)=2^2+3\cdot2-1=9.
\]
\end{examplebox}

For polynomials, this direct substitution works because polynomial functions are continuous everywhere. Rational functions behave similarly at points where the denominator is not zero.

\subsection*{When direct substitution fails}

Direct substitution may produce expressions such as
\[
\frac00,
\qquad
\frac{\infty}{\infty},
\qquad
0\cdot\infty,
\qquad
\infty-\infty.
\]
These are called indeterminate forms. They are not answers. They indicate that the expression needs to be transformed or analyzed more carefully.

\begin{remarkbox}
An indeterminate form means that competing behaviors are present. The expression may have a finite limit, an infinite limit, or no limit. The form itself does not decide the answer.
\end{remarkbox}

\subsection*{Factoring}

Factoring is often used when direct substitution gives \(0/0\).

\begin{examplebox}
Compute
\[
\lim_{x\to3}\frac{x^2-4x+3}{2x-6}.
\]
Substitution gives \(0/0\). Factor:
\[
x^2-4x+3=(x-1)(x-3),
\qquad
2x-6=2(x-3).
\]
For \(x\ne3\),
\[
\frac{x^2-4x+3}{2x-6}=\frac{(x-1)(x-3)}{2(x-3)}=\frac{x-1}{2}.
\]
Therefore
\[
\lim_{x\to3}\frac{x^2-4x+3}{2x-6}=\frac{3-1}{2}=1.
\]
\end{examplebox}

The cancellation is valid for \(x\ne3\). This is enough for the limit, because the limit studies values near \(3\), not necessarily at \(3\).

\subsection*{Rationalizing}

Expressions with radicals may require multiplication by the conjugate.

\begin{examplebox}
Compute
\[
\lim_{x\to25}\frac{\sqrt{x}-5}{x-25}.
\]
Substitution gives \(0/0\). Multiply by the conjugate:
\[
\frac{\sqrt{x}-5}{x-25}\cdot\frac{\sqrt{x}+5}{\sqrt{x}+5}
=
\frac{x-25}{(x-25)(\sqrt{x}+5)}.
\]
For \(x\ne25\), this becomes
\[
\frac{1}{\sqrt{x}+5}.
\]
Therefore
\[
\lim_{x\to25}\frac{\sqrt{x}-5}{x-25}=\frac1{5+5}=\frac1{10}.
\]
\end{examplebox}

Rationalizing is not a trick for its own sake. It removes a difference of square roots that hides the limiting behavior.

\subsection*{Rewriting before deciding}

A common mistake is to see \(0/0\) and conclude that the limit is zero or undefined. This is incorrect. The form \(0/0\) tells us that direct substitution has not answered the question.

\begin{summarybox}
When direct substitution gives an indeterminate form, ask:
\begin{itemize}
\item Can I factor the numerator or denominator?
\item Is there a common factor that cancels for nearby inputs?
\item Is there a square root expression that can be rationalized?
\item Can I divide by a dominant power for limits at infinity?
\item After rewriting, what expression describes the same nearby behavior?
\end{itemize}
\end{summarybox}

The algebra of limits is useful because it helps us reveal behavior. The algebra should serve the limit, not replace the meaning of the limit.